\documentclass{article}
\usepackage[UTF8]{ctex}
\usepackage{amsmath}
\usepackage{amssymb}
\usepackage{graphicx}
\usepackage{geometry}
\geometry{a4paper, margin=2cm}

\newcommand{\problemname}[1]{\begin{center}\Large\bfseries #1\end{center}}
\usepackage{enumitem}

\begin{document}
\problemname{蒜头君买文具}

蒜头君作为学习委员要为班级里的小伙伴们买文具，已知蒜头君的班级里一共有 $n$ 人，其中一个男生需要 $m$ 支铅笔，一个女生需要 $2\timesm$ 支铅笔, 蒜头君一共买了 $k$ 支铅笔，请问蒜头君班级里有多少女生？

\subsection*{输入格式}
输入有三行，第一行是一个正整数 $n(m<n\le100)$，表示班级总人数；

第二行是一个正整数 $m(1\lem\le5)$，表示一个男生所需的铅笔数；

第三行是一个正整数 $k(m\timesn<k\le1000)$，表示蒜头君一共买的铅笔数目。

\subsection*{输出格式}
输出一个整数，表示班级中女生的人数。

给出的整数 $k$ 可以满足男生和女生的比例关系。

\subsection*{样例输入}
\begin{verbatim}
50
2
120
\end{verbatim}

\subsection*{样例输出}
\begin{verbatim}
10
\end{verbatim}


\end{document}
